\chapter{Apprendre sur des covariances: quand le critère sort du modèle}
\label{ch:learning}

\section{Introduction}
\label{sec:learning-introduction}

\begin{contributions}
  Ce chapitre rend compte des travaux suivants, dont les versions publiées sont
  reproduites en fin de chapitre:
  \begin{itemize}
    \item \fullcite{ECML23};
    \item \fullcite{ICML24}.
  \end{itemize}
\end{contributions}

% ============================ TODO texte (~2 paragraphes) ===========================
% Presentation des deux collaborations. C'est le seul endroit du chapitre ou
% elles apparaissent, et aucune des deux n'est un encadrement --- il faut le
% dire, sous peine de laisser croire le contraire par simple voisinage avec les
% \cref{ch:robust-detection,ch:spdnet}:
%  - \fullcite{ECML23}: collaboration avec Alexandre Hippert-Ferrer (alors
%    post-doctorant, aujourd'hui MCF a l'Universite Gustave Eiffel), Florent
%    Bouchard (CR CNRS, L2S), Titouan Vayer (Inria, ENS de Lyon) et Arnaud
%    Breloy (alors MCF a Paris-Nanterre, aujourd'hui professeur au CNAM).
%    PRECISER la nature exacte de la contribution personnelle: ce chapitre est
%    le seul du memoire ou l'auteur n'est ni le premier auteur ni l'encadrant,
%    autant l'assumer clairement.
%  - \fullcite{ICML24}: collaboration avec Florent Bouchard, Malik Tiomoko
%    (alors Huawei Paris), Guillaume Ginolhac (LISTIC) et Frederic Pascal (L2S).
%    C'est la seule publication ICML du dossier; le dire, et dire ce qu'elle a
%    ouvert (la piste rmt reapparait au \cref{ch:perspectives}).
% ====================================================================================

Le chapitre qui précède s'est déroulé dans une situation confortable, et il
faut commencer par dire en quoi. Le \cref{ch:robust-detection} construit des
détecteurs à partir d'un modèle spécifié jusqu'au bout, et c'est de ce modèle
que sort la réponse: le rapport de vraisemblance, le test de Rao et le test de
Wald y sont des \emph{dérivations}, non des inventions. Mais le confort ne
s'arrête pas là, et c'est le point que l'on a peut-être lu trop vite: les
critères qui jugent ces détecteurs --- la constance du taux de fausse alarme,
la courbe $\pdet$ contre $\snr$ --- sont eux aussi internes au modèle. Celui-ci
fournit donc deux choses à la fois: la réponse, et de quoi la noter.

Ce chapitre est celui où la seconde de ces deux fonctions disparaît, et où la
première doit s'en accommoder. Ce n'est pas seulement affaire de dimension. On dit souvent que
l'apprentissage sur des matrices de covariance se distingue par son régime
d'échantillonnage --- peu d'observations pour beaucoup de variables --- et
c'est exact, mais c'est un \emph{mécanisme}, pas une raison. La raison est
ailleurs: dans les deux travaux dont ce chapitre rend compte, ce que l'on
cherche n'est pas une décision ni un paramètre, mais un objet que le modèle ne
sait pas noter. Dans le premier, c'est un graphe de dépendances conditionnelles
que l'on veut pouvoir \emph{lire}; dans le second, une partition d'une
collection de covariances que l'on veut voir coïncider avec une vérité terrain.
Aucune vraisemblance ne mesure la lisibilité d'un graphe, aucune probabilité de
fausse alarme ne mesure la qualité d'une partition. Le critère a quitté le
modèle.

Il ne s'ensuit pas que le modèle disparaisse: c'est là toute la thèse du
chapitre, et elle tient en une phrase. \emph{Faire de l'apprentissage sur des
covariances ne dispense pas du modèle; celui-ci revient sous forme de
correction.} Il revient de deux façons, et les deux travaux illustrent chacune
la sienne. Dans la \cref{sec:learning-graphes}, il revient comme
\textbf{contrainte}: la vraisemblance elliptique, pénalisée pour favoriser la
parcimonie et contrainte à un modèle factoriel de rang faible, fournit
l'estimateur, et c'est la structure imposée qui rend le graphe lisible. Dans la
\cref{sec:learning-frechet}, il revient comme \textbf{correction asymptotique}:
l'estimateur naïf --- la moyenne de Fréchet des covariances empiriques --- est
biaisé dans le régime où l'on travaille, et la théorie des matrices aléatoires
fournit le terme correctif qui le rend consistant. Contrainte d'un côté,
correction de l'autre; dans les deux cas, un objet statistique qui n'est plus
là pour répondre mais pour rattraper.

Le mouvement ne s'arrête pas là, et autant l'annoncer tout de suite pour que
ce chapitre se lise à sa place. Le \cref{ch:spdnet}, qui suit, ira jusqu'au
terme: il n'y subsistera aucune vraisemblance, la réponse viendra des données
par l'intermédiaire de paramètres appris, et le critère y sera une erreur de
classification. Le présent chapitre est le pas intermédiaire, celui où le
modèle a déjà cessé de juger mais n'a pas encore cessé d'estimer. On verra en
temps voulu qu'il ne disparaît pas non plus au chapitre suivant --- il y
revient par la fenêtre, sous la forme du choix d'une moyenne --- mais c'est
l'affaire de ce chapitre-là.

Un mot sur ce que ce chapitre ne fait pas. Les deux travaux qu'il rassemble
sont des articles de conférence, denses et complets, et ils sont reproduits à
la fin du chapitre. Le texte qui suit ne les paraphrase pas: il pose les
modèles, énonce les résultats principaux et discute ce qui se passe quand on
les met en \oe uvre, mais il renvoie systématiquement aux articles pour les
démonstrations, les variantes et le détail des campagnes numériques. C'est la
règle retenue dans toute la \cref{pt:contributions}, et elle se justifie ici
plus qu'ailleurs: le lecteur a les articles sous la main, à
\paperpage{ECML23} et \paperpage{ICML24}.

\paragraph{Ce qu'il faut avoir en tête.} Ce chapitre se lit sans revenir à la
\cref{pt:context}, mais il s'appuie sur quatre choses qui y sont construites.
La covariance empirique et son comportement quand $\Nobs$ n'écrase plus
$\Nvar$, en \cref{sec:context-covariance-regularized-estimation}: c'est le
point de départ des deux sections. Les distributions elliptiques et la
$M$-estimation, en \cref{sec:context-distributions-reel}, qui portent toute la
modélisation de la \cref{sec:learning-graphes}. Le cône $\SPD{\Nvar}$, sa métrique
affine invariante et la distance associée
(\cref{sec:context-riemann-spd,prop:context-riemann-distspd}), ainsi que la
moyenne de Fréchet (\cref{def:context-riemann-frechet}), qui sont l'objet même
de la \cref{sec:learning-frechet}. Enfin l'optimisation riemannienne
(\cref{sec:context-riemann-optimization}), gradient, rétraction et transport,
qui est l'outil commun aux deux travaux: c'est même la seule chose qu'ils
partagent au niveau algorithmique.

Deux notions nouvelles n'ont en revanche pas de place dans l'état des lieux et
sont introduites ici, au strict nécessaire: les modèles graphiques gaussiens
au début de la \cref{sec:learning-graphes}, et la théorie des matrices aléatoires
en \cref{subsec:learning-rmt}. Ce choix est délibéré. Aucune des deux n'est
mobilisée ailleurs dans le mémoire, et la \cref{pt:context} n'a vocation à
contenir que ce qui sert à plusieurs chapitres.

\section{Cadre commun: quand le critère sort du modèle}
\label{sec:learning-cadre}

Cette section n'est dans aucun des deux articles. Elle est écrite pour ce
mémoire, et c'est elle qui fait du chapitre autre chose qu'une juxtaposition:
elle dit ce que les deux travaux ont en commun, et pourquoi ce commun est un
énoncé de modélisation et non une coïncidence de vocabulaire.

\subsection{Deux tâches, un même déplacement}
\label{subsec:learning-taches}

Le point de départ est le même que partout dans ce mémoire: on dispose
d'observations multivariées, et on résume leur structure de second ordre par
une matrice de covariance. Ce qui change est ce que l'on en fait.

\begin{definition}[{Apprentissage d'un graphe de dépendances}]
\label{def:learning-tache-graphe}
  Étant données $\Nobs$ observations $\{\x_\iobs\}_{\iobs=1}^{\Nobs}$ de
  $\R^{\Nvar}$, on cherche un graphe non orienté
  $\graph=(\vertexset,\edgeset)$ à $\Nvar$ sommets, un par variable, dont les
  arêtes traduisent les dépendances conditionnelles entre variables.
\end{definition}

\begin{definition}[{Partitionnement d'une collection de covariances}]
\label{def:learning-tache-partition}
  Étant données $\Nmat$ collections d'observations
  $\{\X_\imat\}_{\imat=1}^{\Nmat}$, chacune de $\Nobs$ vecteurs de
  $\R^{\Nvar}$ et de covariance propre $\MCov_\imat$, on cherche une
  affectation $\{\classlabel_\imat\}_{\imat=1}^{\Nmat}$ des collections à
  $\Nclass$ groupes, qui regroupe les covariances semblables.
\end{definition}

Ces deux problèmes n'ont, en apparence, rien à voir: l'un travaille sur une
seule covariance et cherche à en lire la structure interne, l'autre sur une
collection entière et cherche à l'organiser. Ils partagent pourtant un trait
qui les sépare tous deux du \cref{ch:robust-detection}, et ce trait commande
tout ce qui suit.

\begin{remark}[{Le critère n'est plus dans le modèle}]
\label{rem:learning-critere}
  Dans les deux archétypes de test du \cref{ch:robust-detection}
  (\cref{def:detection-archetype-cible,def:detection-archetype-egalite}), le
  modèle statistique remplit deux fonctions simultanément: il fournit
  l'estimateur --- par maximum de vraisemblance ou $M$-estimation --- et il
  fournit le critère qui juge la décision, à savoir la loi de la statistique
  sous $\hypothesis_0$, dont dérivent le seuil, la $\pfa$ et la propriété cfar.
  Que la dispersion y soit un paramètre de nuisance ou le paramètre testé ne
  change rien à cette double fonction; c'est même l'unique point commun des
  deux archétypes.

  Ici, la seconde fonction tombe. La qualité d'un graphe appris se juge sur sa
  \emph{modularité} ou sur son accord avec une interprétation experte; celle
  d'une partition, sur son accord avec une vérité terrain. Ni l'une ni l'autre
  n'est une fonction de la vraisemblance, et aucune ne s'écrit dans le
  paramètre du modèle. L'estimateur, lui, continue d'en venir --- mais rien,
  dans le modèle, ne dit plus s'il fait l'affaire.
\end{remark}

C'est cette dissymétrie qui rend nécessaire, ici et pas ailleurs, un travail
sur l'estimateur lui-même. Tant que le modèle certifie sa propre sortie, il
suffit de le résoudre; dès qu'il ne la certifie plus, il faut se demander
\emph{où} il échoue, et le corriger là. Les deux sections suivantes répondent à
cette question de deux manières différentes. Avant cela, il faut dire
précisément ce qui échoue.

\subsection{Ce que le régime dimensionnel casse}
\label{subsec:learning-rmt}

% Cette sous-section est le contexte minimal de theorie des matrices
% aleatoires. Elle est ecrite ici plutot qu'en \cref{pt:context} parce que rien
% d'autre dans le memoire ne s'en sert. Elle suit l'ordre de la presentation
% HCERES: regime, Marchenko-Pastur, modele spike, puis la consequence sur les
% distances --- qui est la seule chose dont la \cref{sec:learning-frechet} a
% reellement besoin.

La \cref{sec:context-covariance-regularized-estimation} a déjà signalé que la
covariance empirique se dégrade lorsque $\Nvar$ devient comparable à $\Nobs$,
et l'a fait par le conditionnement. Cette lecture est juste mais grossière:
elle dit que l'estimateur devient instable, elle ne dit pas \emph{comment}. La
théorie des matrices aléatoires le dit exactement, et c'est ce qui permet de
corriger plutôt que de régulariser à l'aveugle.

\begin{definition}[{Régime des matrices aléatoires}]
\label{def:learning-regime-rmt}
  On parle de \textbf{régime des matrices aléatoires} lorsque la dimension et
  le nombre d'observations croissent conjointement,
  \begin{equation}
    \Nvar\to\infty,
    \quad
    \Nobs\to\infty,
    \quad
    \frac{\Nvar}{\Nobs}\to\cratio\in(0,1),
    \label{eq:learning-regime-rmt}
  \end{equation}
  par opposition au \textbf{régime asymptotique classique}, où $\Nvar$ reste
  fixe et $\Nobs\to\infty$. Le réel $\cratio$ est appelé \textbf{rapport de
  concentration}.
\end{definition}

\begin{remark}[{Pourquoi ce régime plutôt que l'autre}]
\label{rem:learning-pourquoi-regime}
  Le régime classique fait de la covariance empirique un estimateur consistant,
  et c'est sur lui que reposent les garanties usuelles. Il décrit mal les
  situations de ce chapitre. Dans le partitionnement hyperspectral du
  \cref{subsec:learning-hyperspectral}, la covariance d'un pixel est estimée
  sur une fenêtre spatiale de quelques dizaines d'échantillons pour une
  dizaine de bandes après réduction; dans les données \textsc{eeg}, la
  situation est du même ordre. Le rapport $\Nvar/\Nobs$ n'y est pas
  négligeable, et le traiter comme s'il l'était revient à supposer résolu le
  seul problème réellement présent. La \cref{def:learning-regime-rmt} n'est pas
  une abstraction: c'est le régime de travail effectif.
\end{remark}

\begin{property}[{Loi de Mar\v{c}enko-Pastur}]
\label{prop:learning-marchenko-pastur}
  Soit $\X\in\R^{\Nvar\times\Nobs}$ à entrées indépendantes, centrées et de
  variance unité, et $\SCM=\frac{1}{\Nobs}\X\X^\T$ la covariance empirique
  associée, de sorte que $\MCov=\Id{\Nvar}$. Dans le régime
  \cref{eq:learning-regime-rmt}, la distribution empirique des valeurs propres
  de $\SCM$ converge vers la densité \sidecite[1.2cm]{marchenko1967distribution}
  \begin{equation}
    \MPdens(\eig)
    =
    \frac{1}{2\pi\cratio\eig}
    \sqrt{(\eig_+-\eig)(\eig-\eig_-)}\,
    \mathbb{1}_{[\eig_-,\eig_+]}(\eig),
    \qquad
    \eig_\pm=(1\pm\sqrt{\cratio})^2 .
    \label{eq:learning-marchenko-pastur}
  \end{equation}
\end{property}

\begin{remark}[{Ce que cette loi dit vraiment}]
\label{rem:learning-lecture-mp}
  L'énoncé mérite d'être lu lentement, car il est plus fort qu'il n'en a
  l'air. La vraie covariance est l'identité: \emph{toutes} ses valeurs propres
  valent $1$. Celles de son estimateur, elles, s'étalent sur tout l'intervalle
  $[(1-\sqrt{\cratio})^2,(1+\sqrt{\cratio})^2]$, et cet étalement ne disparaît
  pas quand on ajoute des données --- il ne dépend que de $\cratio$. À
  $\cratio=1/2$, le spectre observé couvre $[0{,}09;\,2{,}91]$ là où la vérité
  est un point; à $\cratio=1$, il couvre $[0,4]$. Le biais n'est donc ni du
  bruit ni un artefact de petit échantillon: c'est un phénomène déterministe,
  parfaitement décrit, et donc parfaitement corrigible. C'est ce dernier point
  qui fait toute la différence avec la régularisation par rétrécissement de la
  \cref{sec:context-covariance-regularized-estimation}, laquelle contracte le
  spectre sans savoir de combien.
\end{remark}

\wipfig{FIGURE 1/3 --- Marchenko-Pastur}{Histogramme des valeurs propres de la
scm pour $\MCov=\Id{\Nvar}$, superposé à la densité
\cref{eq:learning-marchenko-pastur}, pour trois valeurs de $\cratio$
($0{,}1$, $0{,}5$, $1$). C'est la figure explicative dont dépend toute la remarque
\emph{Ce que cette loi dit vraiment}. Le code existe déjà: dépôt
\texttt{marcenko} des slides HCERES (\texttt{visualize\_standard\_marcenko.py}),
qui exporte déjà en PGFPlots. À porter sous \texttt{qanat} avec les conventions
\texttt{hdrlib} (backend-free, \texttt{make\_mc\_parser}) pour obtenir la ligne
de provenance et le QR code.}

% ============================== FIGURE 1/3 =========================================
% Placeholder. Une fois l'experience portee sous qanat, remplacer par:
%
% \begin{figure}[htb]
%   \centering
%   \input{gfx/generated/learning_marchenko/marchenko/plot.tex}
%   \provMarchenkoQR[1cm]
%   \caption{...\provMarchenko}
%   \label{fig:learning-marchenko}
% \end{figure}
%
% et lancer `just figures` pour que \provMarchenko soit definie.
% ====================================================================================

Le modèle blanc n'est évidemment pas celui des applications: les covariances
que l'on manipule ont une structure. Le modèle de référence est alors le
\textbf{modèle à pics} (\emph{spiked model}), qui superpose quelques directions
énergétiques à un fond isotrope.

\begin{property}[{Déplacement des pics}]
\label{prop:learning-spiked}
  Soit $\MCov=\sum_{\ivar=1}^{\Nrank}\eig_\ivar\vector{u}_\ivar\vector{u}_\ivar^\T+\Id{\Nvar}$
  avec $\eig_\ivar>1$ et $\Nrank$ fixe. Dans le régime
  \cref{eq:learning-regime-rmt}, les valeurs propres de $\SCM$ associées aux
  pics émergent au-dessus du support de
  \cref{eq:learning-marchenko-pastur} et se placent, asymptotiquement, en
  \sidecite[1.2cm]{baik2006eigenvalues,johnstone2001distribution}
  \begin{equation}
    \widehat{\eig}_\ivar
    \;\simeq\;
    \eig_\ivar\left(1+\frac{\cratio}{\eig_\ivar-1}\right),
    \label{eq:learning-spiked}
  \end{equation}
  pourvu que $\eig_\ivar-1>\sqrt{\cratio}$; en deçà de ce seuil, le pic reste
  noyé dans le fond et n'est plus discernable.
\end{property}

Deux enseignements en découlent, et ce sont eux qui portent toute la
\cref{sec:learning-frechet}. D'abord, le biais est \emph{explicite}: la
\cref{eq:learning-spiked} donne le déplacement en fonction de $\cratio$ et de
la valeur propre vraie, donc on peut l'inverser. Ensuite, ce biais affecte tout
ce qui se calcule à partir du spectre --- et la distance affine invariante de
la \cref{prop:context-riemann-distspd} en fait précisément partie, puisqu'elle
ne dépend que des valeurs propres généralisées du couple.

\begin{property}[{Non-consistance de la distance empirique}]
\label{prop:learning-distance-biaisee}
  Soient $\MCov$ une covariance et $\SCM$ son estimateur empirique sur $\Nobs$
  observations. Dans le régime \cref{eq:learning-regime-rmt}, et pour toute
  matrice de référence $\MCov_0\in\SPD{\Nvar}$,
  \begin{equation}
    \geodist^2(\MCov_0,\SCM)
    \;\overset{\text{\footnotesize/}}{\longrightarrow}\;
    \geodist^2(\MCov_0,\MCov).
    \label{eq:learning-distance-biaisee}
  \end{equation}
  Un estimateur consistant $\distrmt^2(\MCov_0,\SCM)$ de $\geodist^2(\MCov_0,\MCov)$
  peut en revanche être construit à partir des seules quantités observables
  $\cratio$ et $\{\eig_\ivar(\MCov_0^{-1}\SCM)\}$
  \sidecite[1.5cm]{couillet2019distances}.
\end{property}

Voilà le point exact où le modèle cesse de suffire, et voilà aussi la forme que
prend son retour. Il ne s'agit pas de renoncer à la distance riemannienne, ni
de la remplacer par une heuristique: il s'agit d'en estimer la valeur vraie à
partir de données bruitées, en utilisant la connaissance précise que l'on a du
bruit. C'est une correction, au sens le plus littéral.

\subsection{Deux façons de faire revenir le modèle}
\label{subsec:learning-deux-facons}

\begin{table}[htb]
  \centering
  \small
  \begin{tabular}{@{}p{0.22\textwidth}p{0.34\textwidth}p{0.34\textwidth}@{}}
    \toprule
    & \textbf{\cref{sec:learning-graphes}} & \textbf{\cref{sec:learning-frechet}} \\
    & apprentissage de graphes & moyenne de Fréchet corrigée \\
    \midrule
    Objet cherché
      & un graphe sur $\Nvar$ sommets
      & un représentant de $\Nmat$ covariances \\
    Critère
      & modularité, lisibilité experte
      & accord avec une vérité terrain \\
    Ce qui échoue
      & $\Nobs\simeq\Nvar$: la scm n'identifie pas de structure, et le modèle
        gaussien est fragile aux queues lourdes
      & régime \cref{eq:learning-regime-rmt}: la distance, donc la moyenne,
        est biaisée \\
    Retour du modèle
      & \textbf{contrainte}: vraisemblance elliptique, pénalisée et contrainte
        à un modèle factoriel
      & \textbf{correction}: terme correctif issu de la théorie des matrices
        aléatoires \\
    Variété
      & $\Factorman{\Nvar}{\Nrank}$, rang faible plus diagonale
      & $\SPD{\Nvar}$ \\
    \bottomrule
  \end{tabular}
  \caption[Les deux façons de faire revenir le modèle]{Les deux sections du
    chapitre en regard. La dernière ligne est la seule chose que les deux
    travaux partagent au niveau algorithmique: dans les deux cas, le problème
    est posé sur une variété et résolu par descente riemannienne
    (\cref{sec:context-riemann-optimization}).}
  \label{tab:learning-deux-facons}
\end{table}

La \cref{tab:learning-deux-facons} met les deux travaux en regard. On y lit
surtout ceci: le déplacement décrit plus haut --- le critère n'est plus dans le
modèle --- ne laisse pas le praticien sans ressource. Il l'oblige à choisir \emph{où} injecter ce
qu'il sait. Contraindre l'espace de recherche, ou corriger l'estimateur dans un
espace inchangé --- ce sont deux stratégies distinctes, et il n'est pas
indifférent que les deux passent par la même boîte à outils riemannienne. C'est
d'ailleurs la raison pour laquelle ces deux travaux, menés avec des équipes
différentes et publiés dans des lieux différents, se lisent ensemble.

\section{Apprendre un graphe: vraisemblance elliptique pénalisée et modèle factoriel}
\label{sec:learning-graphes}

\begin{contributions}
  Cette section reprend \fullcite{ECML23}, reproduit p.~\paperpage{ECML23}.
\end{contributions}

\subsection{Du graphe à la matrice de précision}
\label{subsec:learning-graphe-precision}

Le pont entre un graphe et une matrice est classique, et il passe par
l'inverse de la covariance.

\begin{definition}[{Modèle graphique gaussien}]
\label{def:learning-ggm}
  Soit $\x=[x_1,\dots,x_{\Nvar}]^\T\sim\gaussian(\vector{0},\MCov)$ et
  $\Precision=\inv{\MCov}$ sa \textbf{matrice de précision}. Le champ de Markov
  gaussien associé est le graphe non orienté à $\Nvar$ sommets dans lequel une
  arête relie les sommets $\ivar$ et $\ivarbis$ si et seulement si les
  variables $x_\ivar$ et $x_\ivarbis$ sont conditionnellement dépendantes
  sachant toutes les autres \sidecite[1.0cm]{dempster1972covariance,lauritzen1996graphical}.
\end{definition}

\begin{property}[{Support de la précision et arêtes}]
\label{prop:learning-support}
  Avec les notations de la \cref{def:learning-ggm},
  \begin{equation}
    \mathrm{corr}\!\left(x_\ivar x_\ivarbis \mid x_{\intrange{1}{\Nvar}\setminus\{\ivar,\ivarbis\}}\right)
    =
    -\frac{[\Precision]_{\ivar\ivarbis}}{\sqrt{[\Precision]_{\ivar\ivar}[\Precision]_{\ivarbis\ivarbis}}} .
    \label{eq:learning-corr-partielle}
  \end{equation}
  Une entrée non nulle de $\Precision$ hors diagonale équivaut donc à une arête
  du graphe: \emph{apprendre le graphe, c'est estimer le support de la matrice
  de précision}.
\end{property}

L'approche standard consiste alors à estimer $\Precision$ par maximum de
vraisemblance régularisé,
\begin{equation}
  \max_{\Precision\in\SPD{\Nvar}}
  \;
  \log\Det(\Precision)-\trace{\SCM\Precision}-\regparam\,\penfun(\Precision),
  \label{eq:learning-glasso}
\end{equation}
où $\penfun$ est une pénalité favorisant la parcimonie et $\regparam\in\R_+$ un
paramètre de régularisation. Le choix $\penfun=\norm{\cdot}_1$ donne
l'algorithme \algname{glasso} \sidecite[0cm]{friedman2008sparse}, qui est la
référence du domaine.

\subsection{Deux insuffisances}
\label{subsec:learning-insuffisances}

La formulation \cref{eq:learning-glasso} appelle deux objections, et ce sont
elles qui motivent tout le travail.

\paragraph{L'hypothèse gaussienne.} La vraisemblance de
\cref{eq:learning-glasso} est celle du modèle gaussien, dont la fragilité aux
queues lourdes et aux points aberrants est le fil conducteur du
\cref{ch:context-distributions}. Les données réelles pour lesquelles on
souhaite un graphe --- réponses catégorielles, séries financières, séries
géodésiques --- ne sont pas gaussiennes, et une seule observation atypique
suffit à déplacer des arêtes. Le remède est celui déjà employé au
\cref{ch:robust-detection}: passer à la famille elliptique
(\cref{def:context-distributions-elliptical}), dont la log-vraisemblance
négative s'écrit
\begin{equation}
  \loss(\MCov)
  \;\propto\;
  \frac{1}{\Nobs}\sum_{\iobs=1}^{\Nobs}
  \rho\!\left(\x_\iobs^\T\inv{\MCov}\x_\iobs\right)
  +\frac{1}{2}\log\Det(\MCov),
  \qquad
  \rho(t)=-\log\densgen(t),
  \label{eq:learning-nll-elliptique}
\end{equation}
et qui redonne le cas gaussien pour $\densgen(t)=\exp(-t/2)$
\sidecite[0cm]{vogel2011elliptical}. Dans les expériences des articles, le
générateur retenu est celui de la $t$ de Student à $\dof=5$ degrés de liberté.

\paragraph{L'absence de structure.} La seconde objection est celle du régime
dimensionnel, sous la forme déjà rencontrée en
\cref{sec:context-covariance-regularized-estimation}: lorsque
$\Nobs\simeq\Nvar$ ou $\Nobs<\Nvar$, la covariance empirique qui alimente
\cref{eq:learning-glasso} n'est plus fiable, et la parcimonie seule ne suffit
pas à compenser. Il faut réduire le nombre de paramètres, et la façon
canonique de le faire est de supposer les données engendrées par un modèle
linéaire de faible dimension.

\begin{definition}[{Modèle factoriel de rang faible}]
\label{def:learning-facteur}
  On dit que $\x$ suit un modèle factoriel de rang $\Nrank$ lorsque
  $\x=\Factormat\vector{s}+\vector{\epsilon}$, avec
  $\Factormat\in\R^{\Nvar\times\Nrank}$ et $\vector{s}$, $\vector{\epsilon}$
  indépendants. La covariance est alors de la forme
  $\MCov=\Lowrank+\Diagpart$, où
  $\Lowrank=\Factormat\,\E[\vector{s}\vector{s}^\T]\Factormat^\T$ appartient à
  l'ensemble $\PSDrank{\Nvar}{\Nrank}$ des matrices semi-définies positives de
  rang $\Nrank$, et $\Diagpart\in\Diagpos{\Nvar}$. On note
  \begin{equation}
    \Factorman{\Nvar}{\Nrank}
    =
    \left\{\MCov=\Lowrank+\Diagpart,\;
    \Lowrank\in\PSDrank{\Nvar}{\Nrank},\;
    \Diagpart\in\Diagpos{\Nvar}\right\}
    \label{eq:learning-factorman}
  \end{equation}
  l'ensemble des matrices « rang faible plus diagonale ».
\end{definition}

\begin{remark}[{Le gain en paramètres}]
\label{rem:learning-gain-parametres}
  Cette paramétrisation fait passer la dimension du problème d'estimation de
  $\Nvar(\Nvar+1)/2$ à $\Nvar(\Nrank+1)-\Nrank(\Nrank-3)/2$. Pour
  $\Nvar=33$ et $\Nrank=10$ --- les valeurs du jeu \texttt{animals} de la
  \cref{subsec:learning-graphe-resultats} --- cela fait passer de $561$ à
  $328$ paramètres. Le gain est réel sans être spectaculaire; l'intérêt du
  modèle factoriel n'est d'ailleurs pas seulement là. La
  \cref{prop:learning-composantes} donne la vraie raison.
\end{remark}

\begin{property}[{Rang et composantes connexes}]
\label{prop:learning-composantes}
  Dans les modèles à contrainte laplacienne, une matrice de précision de rang
  $\Nrank$ correspond à un graphe à $\Nvar-\Nrank$ composantes connexes
  \sidecite[1.5cm]{kumar2020unified}.
\end{property}

C'est ce dernier énoncé qui transforme une contrainte de calcul en une
contrainte \emph{interprétable}: la structure spectrale que l'on impose pour
des raisons d'estimation gouverne aussi le nombre de groupes que le graphe
appris fera apparaître. Là où les méthodes concurrentes exigent que ce nombre
soit fixé a priori, il découle ici du rang.

\subsection{Le problème posé, et pourquoi il est difficile}
\label{subsec:learning-probleme-graphe}

En combinant les deux remèdes, le problème d'apprentissage s'écrit
\begin{equation}
  \min_{\MCov\in\SPD{\Nvar}}
  \;
  \loss(\MCov)+\regparam\,\penfun(\MCov)
  \qquad\text{sous contrainte}\qquad
  \MCov\in\Factorman{\Nvar}{\Nrank},
  \label{eq:learning-probleme-gfm}
\end{equation}
avec $\loss$ la vraisemblance elliptique \cref{eq:learning-nll-elliptique} et,
pour pénalité, une fonction élément par élément appliquée à la
\emph{précision}:
\begin{equation}
  \penfun(\MCov)=\sum_{\ivar\neq\ivarbis}\penaltyelem\!\left([\inv{\MCov}]_{\ivar\ivarbis}\right).
  \label{eq:learning-penalite}
\end{equation}

\begin{remark}[{La difficulté est dans le mélange}]
\label{rem:learning-difficulte}
  Aucun des deux ingrédients n'est neuf pris isolément: les modèles graphiques
  elliptiques existent, les modèles factoriels aussi. Ce qui résiste, c'est
  leur réunion, et pour une raison précise: la vraisemblance
  \cref{eq:learning-nll-elliptique} et la contrainte
  \cref{eq:learning-factorman} portent sur la \emph{covariance}, tandis que la
  pénalité \cref{eq:learning-penalite} porte sur son \emph{inverse}. Le
  problème mêle donc des termes exprimés dans les deux représentations, ce
  qu'aucun schéma de descente coordonnée à la \algname{glasso} ne sait
  traiter. La réponse habituelle est le fractionnement de variables et
  l'\algname{admm}, qui résout un problème approché et introduit ses propres
  paramètres de réglage.
\end{remark}

\subsection{Résolution par optimisation riemannienne}
\label{subsec:learning-resolution-graphe}

La contribution de \fullcite{ECML23} est de traiter
\cref{eq:learning-probleme-gfm} directement, en tirant parti du fait que
$\Factorman{\Nvar}{\Nrank}$ est une sous-variété lisse de $\SPD{\Nvar}$. Le
problème contraint devient alors un problème \emph{libre} sur une variété, et
la machinerie de la \cref{sec:context-riemann-optimization} s'applique sans
approximation ni paramètre supplémentaire.

Sur $\SPD{\Nvar}$ --- c'est-à-dire quand on relâche la contrainte de rang ---
les trois ingrédients sont ceux du \cref{ch:context-riemann}: la métrique
affine invariante \cref{eq:context-riemann-metriqueaffine}, le transport
parallèle \cref{eq:context-riemann-transportspd}, et la rétraction du second
ordre
\begin{equation}
  \Retr_{\MCov}(\tvec)=\MCov+\tvec+\tfrac{1}{2}\tvec\inv{\MCov}\tvec,
  \label{eq:learning-retraction-spd}
\end{equation}
qui garantit la positivité des itérées. Le gradient riemannien s'obtient du
gradient euclidien par
$\Rgrad f(\MCov)=\MCov\,\symop(\nabla f(\MCov))\,\MCov$.

\begin{proposition}[{Gradient riemannien de la vraisemblance elliptique}]
\label{prop:learning-gradient-nll}
  Le gradient riemannien de $\loss$ en $\MCov\in\SPD{\Nvar}$ vaut
  \begin{equation}
    \Rgrad\loss(\MCov)
    =
    \frac{1}{2}\MCov
    -\frac{1}{2\Nobs}\sum_{\iobs=1}^{\Nobs}
    \weight\!\left(\x_\iobs^\T\inv{\MCov}\x_\iobs\right)\x_\iobs\x_\iobs^\T,
    \qquad
    \weight(t)=-2\frac{\densgen'(t)}{\densgen(t)},
    \label{eq:learning-gradient-nll}
  \end{equation}
  et celui de la pénalité \cref{eq:learning-penalite} a pour entrées
  $[\Rgrad\penfun(\MCov)]_{\ivar\ivarbis}=\penaltyelem'([\inv{\MCov}]_{\ivar\ivarbis})$
  hors diagonale, et $0$ sur la diagonale.
\end{proposition}

\begin{remark}[{La fonction de poids, encore}]
\label{rem:learning-poids}
  La fonction $\weight$ de \cref{eq:learning-gradient-nll} est exactement celle
  de la $M$-estimation de la \cref{subsec:context-robust-m-estimation}:
  $\weight(t)=1$ dans le cas gaussien, $\weight(t)=(\dof+\Nvar)/(\dof+t)$ pour
  la $t$ de Student. La robustesse ne vient donc d'aucun mécanisme nouveau ---
  elle est la même pondération des observations par leur distance de
  Mahalanobis, que celle qui produit l'estimateur de Tyler au
  \cref{ch:robust-detection}. C'est un exemple net de ce que ce mémoire essaie
  de montrer: le même objet, ici une fonction de poids, sert à décider dans un
  chapitre et à apprendre dans un autre.
\end{remark}

Pour le modèle factoriel, la géométrie retenue est celle des matrices
semi-définies positives de rang fixe
\sidecite[-0.4cm]{bonnabel2009riemannian,bouchard2021lowrank}, combinée à celle du
cône diagonal. Le détail de cette construction --- paramétrisation, projections,
rétraction sur $\Factorman{\Nvar}{\Nrank}$ --- est dans l'article; on retiendra
ici son effet pratique.

\begin{remark}[{Le vrai gain est algorithmique}]
\label{rem:learning-gain-algorithmique}
  Une itération de la méthode proposée ne demande que la \textsc{svd} tronquée
  du facteur de rang faible, là où les méthodes à contrainte laplacienne
  concurrentes demandent la décomposition complète du laplacien ou de la
  matrice d'adjacence à chaque itération. Le goulot d'étranglement passe donc
  de $\bigO(\Nvar^3)$ à $\bigO(\Nvar\Nrank^2)$. Comme au
  \cref{ch:robust-detection}, la structure ne sert pas seulement à mieux
  estimer: elle sert aussi à calculer moins.
\end{remark}

\subsection{Ce que la structure achète}
\label{subsec:learning-graphe-resultats}

% ============================ TODO texte (~2 paragraphes) ===========================
% Discussion des resultats. Le materiau chiffre est dans l'article (fig. 1 et 2)
% et les figures ne sont PAS rejouees ici (cf. la convention annoncee plus bas):
%  - jeu `animals` (Osherson/Lake--Tenenbaum): \Nvar = 33 animaux, \Nobs = 102
%    questions binaires, donc des donnees categorielles franchement non
%    gaussiennes. Modularites relevees dans l'article: glasso 0,56; ngl 0,38;
%    sgl 0,74; stgl 0,84; ggm 0,54; ggfm 0,79; egm 0,44; egfm 0,80. Rang fixe
%    a 10 pour les modeles factoriels, 8 composantes imposees pour sgl/stgl.
%  - LE POINT A MARTELER: ggfm et egfm atteignent la modularite de sgl et stgl
%    SANS qu'on leur donne le nombre de composantes --- celui-ci decoule du
%    rang (\cref{prop:learning-composantes}). C'est exactement la these de la
%    section: la contrainte remplace la supervision.
%  - jeu GNSS (reseau de stations du Piton de la Fournaise, 33 stations):
%    verifier la description exacte dans l'article avant d'ecrire, et dire ce
%    que le graphe appris donne a voir (regroupement geographique des stations
%    du sommet vs. celles de la base).
%  - Ne PAS terminer sur la modularite: c'est un critere de substitution. Dire
%    que le juge final reste la lecture experte du graphe, ce qui illustre
%    exactement la remarque « Le critere n'est plus dans le modele ».
% ====================================================================================

\begin{remark}[{Pourquoi aucune figure ici}]
\label{rem:learning-pas-de-figure}
  Les graphes appris sur les jeux \texttt{animals} et \textsc{gnss} figurent
  dans l'article reproduit, p.~\paperpage{ECML23}, et ne sont pas reproduits
  ici. La règle de reproductibilité annoncée au chapitre
  \hyperref[ch:commentlire]{\emph{Comment lire ce mémoire}} veut que toute
  figure de simulation du mémoire porte sa ligne de provenance et son
  expérience rejouable; recopier une figure d'article sans cette chaîne
  reviendrait à en afficher la forme sans la garantie. Le lecteur intéressé par
  les graphes eux-mêmes les trouvera à quelques pages d'ici, dans leur contexte
  d'origine.
\end{remark}

\section{Corriger la moyenne: matrices aléatoires et moyenne de Fréchet}
\label{sec:learning-frechet}

\begin{contributions}
  Cette section reprend \fullcite{ICML24}, reproduit p.~\paperpage{ICML24}.
\end{contributions}

\subsection{Classer et partitionner avec des moyennes}
\label{subsec:learning-centroides}

Il faut commencer par la tâche, car c'est elle qui fait apparaître l'objet à
corriger. On dispose de $\Nmat$ collections d'observations
$\{\X_\imat\}_{\imat=1}^{\Nmat}$, chacune résumée par sa covariance empirique
$\SCM_\imat=\frac{1}{\Nobs}\X_\imat\X_\imat^\T$, et l'on veut soit leur
attribuer une classe, soit les regrouper. La méthode de référence sur le cône
est d'une simplicité désarmante.

\begin{definition}[{Classification au plus proche centroïde}]
\label{def:learning-centroide}
  Phase d'apprentissage: pour chaque classe $\iclass\in\intrange{1}{\Nclass}$,
  on calcule la moyenne de Fréchet des covariances de cette classe,
  \begin{equation}
    \Fmeanclass{\iclass}
    =
    \argmin_{\MCov\in\SPD{\Nvar}}
    \sum_{\imat\,:\,\classlabel_\imat=\iclass}
    \geodist^2(\MCov,\SCM_\imat).
    \label{eq:learning-centroide-fit}
  \end{equation}
  Phase de décision: une covariance $\SCM$ non étiquetée reçoit l'étiquette
  \begin{equation}
    \classlabel
    =
    \argmin_{\iclass\in\intrange{1}{\Nclass}}
    \;
    \geodist^2(\Fmeanclass{\iclass},\SCM).
    \label{eq:learning-centroide-predict}
  \end{equation}
\end{definition}

Le partitionnement en $\Nclass$ groupes suit le même schéma, à ceci près que
les étiquettes sont inconnues: on alterne une étape d'affectation
\cref{eq:learning-centroide-predict} et une étape de recalcul des centroïdes
\cref{eq:learning-centroide-fit}, exactement comme les $\Nclass$-moyennes
euclidiennes \sidecite[0cm]{barachant2011multiclass,congedo2015approximate}.

\begin{remark}[{Tout repose sur la moyenne}]
\label{rem:learning-tout-repose}
  Ces deux algorithmes n'ont, au fond, qu'un seul ingrédient statistique: la
  moyenne de Fréchet de la \cref{def:context-riemann-frechet}. Le reste est de
  la comparaison de distances. Il s'ensuit que la qualité de la classification
  ou de la partition est entièrement gouvernée par la qualité de cet
  estimateur --- et que si celui-ci est biaisé, tout l'est. C'est là que la
  \cref{prop:learning-distance-biaisee} devient un problème pratique et non une
  curiosité asymptotique.
\end{remark}

\subsection{La distance corrigée}
\label{subsec:learning-distance-corrigee}

\begin{remark}[{Une normalisation à signaler}]
\label{rem:learning-normalisation}
  Les travaux de cette section emploient la distance affine invariante
  \emph{normalisée},
  $\geodist^2(\MCov_0,\MCov_1)=\frac{1}{2\Nvar}\norm{\logm(\MCov_0^{-1/2}\MCov_1\MCov_0^{-1/2})}_\frob^2$,
  soit la \cref{eq:context-riemann-distspd} divisée par $2\Nvar$. La
  normalisation par la dimension est sans effet sur les minimiseurs --- donc
  sur la \cref{def:learning-centroide} --- mais elle est indispensable pour que
  la quantité admette une limite dans le régime
  \cref{eq:learning-regime-rmt}, où $\Nvar\to\infty$. Les formules qui suivent
  la supposent en place.
\end{remark}

\begin{proposition}[{Distance de Fisher corrigée par la théorie des matrices aléatoires}]
\label{prop:learning-distance-corrigee}
  Soient $\MCov_0\in\SPD{\Nvar}$, $\SCM$ la covariance empirique de $\MCov$ sur
  $\Nobs$ observations, et
  $\eigvec=[\eig_1,\dots,\eig_{\Nvar}]^\T$ les valeurs propres de
  $\MCov_0^{-1}\SCM$. Alors
  \begin{multline}
    \distrmt^2(\MCov_0,\SCM)
    =
    \frac{1}{\Nvar}\sum_{\ivar=1}^{\Nvar}
    \left[\tfrac{1}{2}\log^2\eig_\ivar+\log\eig_\ivar\right]
    \\
    -(\eigvec-\rmtaux)^\T
    \left[\frac{1}{\Nvar}\mat{Q}\ones{\Nvar}+\frac{1-\cratio}{\cratio}\vector{q}\right]
    -\frac{1-\cratio}{2\cratio}\log^2(1-\cratio),
    \label{eq:learning-distance-corrigee}
  \end{multline}
  où $\rmtaux$, $\mat{Q}$ et $\vector{q}$ sont des quantités auxiliaires
  entièrement déterminées par $\cratio$ et $\eigvec$, est un estimateur
  consistant de $\geodist^2(\MCov_0,\MCov)$ dans le régime
  \cref{eq:learning-regime-rmt}
  \sidecite[1.5cm]{couillet2019distances}:
  \begin{equation}
    \distrmt^2(\MCov_0,\SCM)
    \;\xrightarrow[\Nvar/\Nobs\to\cratio]{}\;
    \geodist^2(\MCov_0,\MCov).
    \label{eq:learning-consistance}
  \end{equation}
\end{proposition}

\begin{remark}[{Comment lire cette formule}]
\label{rem:learning-lecture-correction}
  Le détail des quantités auxiliaires figure dans l'article et n'a pas besoin
  d'être reproduit ici; ce qui compte est la \emph{forme} de
  \cref{eq:learning-distance-corrigee}. Le premier terme est la distance
  empirique naïve, à un décalage près. Tout le reste est du correctif, et ce
  correctif ne fait intervenir que des quantités observables: le rapport
  $\cratio=\Nvar/\Nobs$, que l'on connaît, et le spectre $\eigvec$, que l'on
  mesure. Rien n'exige de connaître $\MCov$. C'est précisément ce qui rend la
  correction utilisable: elle transforme une connaissance théorique du biais en
  une opération sur les données. On notera aussi que le correctif s'évanouit
  quand $\cratio\to0$, c'est-à-dire quand on retombe dans le régime
  asymptotique classique --- la formule contient l'ancienne comme cas limite.
\end{remark}

\subsection{La moyenne corrigée et son gradient}
\label{subsec:learning-moyenne-corrigee}

L'idée de la contribution tient en une substitution: remplacer $\geodist^2$ par
$\distrmt^2$ dans la définition même de la moyenne de Fréchet.

\begin{definition}[{Moyenne de Fréchet corrigée}]
\label{def:learning-frechet-rmt}
  Étant données $\Nmat$ covariances empiriques
  $\{\SCM_\imat\}_{\imat=1}^{\Nmat}$, on appelle moyenne de Fréchet corrigée
  tout minimiseur de
  \begin{equation}
    \Fmeanrmt
    \in
    \argmin_{\MCov\in\SPD{\Nvar}}
    \;
    \frac{1}{\Nmat}\sum_{\imat=1}^{\Nmat}
    \distrmt^2(\MCov,\SCM_\imat).
    \label{eq:learning-frechet-rmt}
  \end{equation}
\end{definition}

\begin{remark}[{Une substitution qui coûte cher}]
\label{rem:learning-substitution-couteuse}
  La \cref{def:learning-frechet-rmt} est immédiate à écrire et beaucoup moins à
  résoudre. La moyenne de Fréchet ordinaire dispose, pour la métrique affine
  invariante, d'une condition d'optimalité de forme close
  \cref{eq:context-riemann-frechetopt} qui donne directement l'algorithme de
  point fixe usuel; rien de tel ici, car $\distrmt^2$ n'est pas la distance
  d'une métrique riemannienne --- c'est un estimateur d'une distance. La
  \cref{eq:learning-frechet-rmt} reste néanmoins un problème d'optimisation sur
  $\SPD{\Nvar}$, et c'est la voie retenue: descente riemannienne avec la
  rétraction \cref{eq:learning-retraction-spd}. Ce qui manque alors est le
  gradient, et c'est lui la contribution technique de l'article.
\end{remark}

Le gradient s'obtient en deux temps. D'abord, on remarque que le coût ne dépend
de $\MCov$ qu'à travers le spectre de $\MCov^{-1/2}\SCM_\imat\MCov^{-1/2}$, ce
qui permet de ramener le calcul sur le cône diagonal.

\begin{proposition}[{Réduction au spectre}]
\label{prop:learning-gradient-reduction}
  Soit $f_\imat(\MCov)=g(\Eigmat_\imat)$ où
  $\MCov^{-1/2}\SCM_\imat\MCov^{-1/2}=\Unit_\imat\Eigmat_\imat\Unit_\imat^\T$
  est une décomposition spectrale. Alors
  \begin{equation}
    \Rgrad f_\imat(\MCov)
    =
    -\MCov^{1/2}\Unit_\imat\,
    \inv{\Eigmat_\imat}\,\Rgrad g(\Eigmat_\imat)\,
    \Unit_\imat^\T\MCov^{1/2},
    \label{eq:learning-gradient-reduction}
  \end{equation}
  où $\Rgrad g(\Eigmat_\imat)$ est le gradient riemannien de $g$ en
  $\Eigmat_\imat$ dans $\Diagpos{\Nvar}$.
\end{proposition}

Il reste à calculer $\Rgrad g$, c'est-à-dire à dériver
\cref{eq:learning-distance-corrigee} par rapport aux valeurs propres. Le
résultat est une expression longue mais entièrement close, faisant intervenir
trois matrices $\mat{A}$, $\mat{B}$, $\mat{C}$ dont les entrées sont des
fonctions explicites des $\eig_\ivar$ et de $\Nobs$; on la trouvera à
p.~\paperpage{ICML24}. Le gradient du coût complet en découle par moyenne,
$\Rgrad h(\MCov)=\frac{1}{\Nmat}\sum_\imat\Rgrad f_\imat(\MCov)$, et
l'\cref{alg:learning-frechet-rmt} s'écrit alors sans difficulté.

\begin{algorithm}[htb]
  \caption{Moyenne de Fréchet corrigée par la théorie des matrices aléatoires}
  \label{alg:learning-frechet-rmt}
  \begin{algorithmic}[1]
    \Require $\{\X_\imat\}_{\imat=1}^{\Nmat}$ dans $\R^{\Nvar\times\Nobs}$,
      point initial $\MCov_0$, tolérance $\varepsilon>0$, budget $\Niter$
    \For{$\imat\in\intrange{1}{\Nmat}$}
      \State $\SCM_\imat\gets\frac{1}{\Nobs}\X_\imat\X_\imat^\T$
    \EndFor
    \Repeat
      \State calculer $\Rgrad h(\MCov_\iiter)$ par
        \cref{eq:learning-gradient-reduction}
      \State calculer le pas $t_\iiter$ par recherche linéaire
      \State $\MCov_{\iiter+1}\gets\Retr_{\MCov_\iiter}\!\left(-t_\iiter\Rgrad h(\MCov_\iiter)\right)$
      \State $\iiter\gets\iiter+1$
    \Until{$\geodist^2(\MCov_\iiter,\MCov_{\iiter-1})<\varepsilon$ ou $\iiter>\Niter$}
    \Ensure $\Fmeanrmt=\MCov_\iiter$
  \end{algorithmic}
\end{algorithm}

\begin{remark}[{Dériver à la main, encore}]
\label{rem:learning-deriver-main}
  Le lecteur qui poursuivra au \cref{ch:spdnet} y retrouvera le même geste, et
  le même argument: l'objet géométriquement juste n'a pas de forme close, sa
  minimisation est itérative, et le travail consiste à en dériver le gradient
  explicitement plutôt qu'à laisser un outil de différentiation automatique
  parcourir les itérations. La raison n'y sera pourtant pas la même qu'ici. Là-bas,
  il s'agira d'économiser la mémoire d'un graphe de calcul qui stockerait
  sinon toutes les itérées; ici, plus simplement, d'avoir un algorithme du
  tout. Le constat, lui, est identique --- et c'est un des rares fils qui
  traverse les deux chapitres au niveau technique et non au niveau des idées.
\end{remark}

\subsection{Ce que la correction achète}
\label{subsec:learning-mse}

Le premier critère est interne: à quelle distance de la vraie moyenne se trouve
la moyenne estimée? C'est une erreur quadratique moyenne, mesurée sur données
synthétiques où la vérité est connue, et c'est le critère qui mesure le
\emph{mécanisme} décrit en \cref{subsec:learning-rmt}.

\wipfig{FIGURE 2/3 --- eqm contre nombre d'échantillons}{eqm (en dB) de
l'estimation de la moyenne de Fréchet en fonction de $\Nobs$, à $\Nvar$ et
$\Nmat$ fixés, pour cinq estimateurs: scm, Ledoit-Wolf linéaire, oas,
Ledoit-Wolf non linéaire, et la moyenne corrigée. Courbe moyenne et bande
interpercentile 5/95. \textbf{Les données existent déjà}: fichiers
\texttt{figures/mse\_nsamples/\{mean,5,95\}.csv} du dépôt de slides HCERES, et
le code est public (\texttt{github.com/AmmarMian/icml-rmt-2024}). Ce qui manque
est le passage sous \texttt{qanat} pour la ligne de provenance, et le
réhabillage aux couleurs du mémoire --- l'export actuel est aux couleurs du
thème L2S, qui n'ont rien à faire dans le \texttt{.tex}.}

\wipfig{FIGURE 3/3 --- eqm contre nombre de matrices}{Même chose en fonction de
$\Nmat$, à $\Nobs$ fixé. C'est la figure qui porte l'argument spécifique de
l'article: le gain \emph{croît} avec le nombre de matrices moyennées, alors que
l'intuition voudrait qu'accumuler des matrices compense le manque
d'échantillons. Données dans \texttt{figures/mse\_nmatrices/}. Contrainte:
partager axes, unités et jeu d'estimateurs avec la figure 2/3.}

% ============================== FIGURES 2/3 et 3/3 =================================
% Placeholders. Une fois les experiences portees sous qanat:
%
% \begin{figure}[htb]
%   \centering
%   \input{gfx/generated/learning_frechet_mse/mse-nsamples/plot.tex}
%   \provMseNsamplesQR[1cm]
%   \caption{...\provMseNsamples}
%   \label{fig:learning-mse-nsamples}
% \end{figure}
%
% idem pour mse-nmatrices. Penser a lancer `just figures`.
% ====================================================================================

% ============================ TODO texte (~2 paragraphes) ===========================
% Commentaire des deux figures, a ecrire une fois les runs qanat faits. Points
% a couvrir, dans cet ordre:
%  1. lecture de la figure 2/3: l'ecart entre la scm et la moyenne corrigee est
%     de plusieurs dizaines de dB a faible \Nobs et se referme quand \Nobs
%     croit --- c'est la signature d'un biais de regime et non d'une variance;
%  2. lecture de la figure 3/3, qui est le point NON TRIVIAL: le gain croit
%     avec \Nmat. L'expliquer: moyenner davantage de matrices reduit la
%     variance mais PAS le biais, qui est commun a toutes les scm; passe un
%     certain \Nmat, c'est le biais qui domine et lui seul distingue les
%     methodes. C'est l'argument qui justifie a lui seul la correction;
%  3. situer les concurrents: Ledoit-Wolf et oas regularisent chaque covariance
%     AVANT de moyenner, la methode proposee corrige la DISTANCE. Dire que ce
%     n'est pas le meme geste.
% ====================================================================================

\subsection{Partitionnement hyperspectral}
\label{subsec:learning-hyperspectral}

L'erreur quadratique est un critère interne, et l'on a dit plus haut --- dans
la remarque \emph{Le critère n'est plus dans le modèle} --- que ce n'était pas
celui qui juge réellement ce chapitre. Le juge est en aval: la partition produite.

% ============================ TODO texte (~2 paragraphes) ===========================
% Resultats de partitionnement. L'article couvre \textsc{eeg} et
% hyperspectral; la section peut se limiter a l'hyperspectral, qui est le
% terrain du memoire, et renvoyer a l'article pour l'\textsc{eeg}.
% Materiau disponible: le run qanat `hs_clustering/run_26` des slides HCERES
% existe deja (Indian Pines, pca_dim dans {5,10,12,16,24}, window_size dans
% {5,7}, n_init=10, seed=1234) et produit les cartes gt / SCM-Kmeans /
% LWF-Kmeans / RMT-Kmeans.
% DECISION A PRENDRE: soit on rejoue ce run sous qanat pour en faire une
% quatrieme figure du memoire (il faut les donnees Indian Pines en local), soit
% on renvoie a l'article. Si on le rejoue, c'est LA figure qui porte la these
% du chapitre, et elle merite alors de fermer la section.
% Points a dire dans tous les cas:
%  - la chaine complete: fenetre spatiale -> scm par pixel -> pca -> K-moyennes
%    riemanniennes; c'est la que \Nvar/\Nobs se lit concretement;
%  - le gain se maintient en aval de l'estimation, donc le biais de la moyenne
%    n'etait pas une subtilite theorique;
%  - refermer sur la remarque « Le critere n'est plus dans le modele »: c'est
%    que le critere qui juge le travail apparait.
% ====================================================================================

\section{Bilan}
\label{sec:learning-bilan}

\wip{Bilan ch. apprentissage: les deux mouvements}{Environ 4 paragraphes.
Nommer les personnes avec leur statut exact (aucun nom n'apparaît dans le
chapitre en dehors du bloc \emph{contributions}), en rappelant que ni l'une ni
l'autre collaboration n'est un encadrement. Dire ce qui a été construit, en
clair et sans équations: d'un côté un cadre d'apprentissage de graphes qui
n'exige plus qu'on lui donne le nombre de composantes, de l'autre une moyenne
de Fréchet qui reste juste quand les covariances sont mal estimées. Pointer les
articles reproduits par leur folio (ECML23 p.~\pageref{paper:ECML23}, ICML24
p.~\pageref{paper:ICML24}). Dire où cela continue --- la piste rmt reparaît au
\cref{ch:perspectives}.}

\wip{Bilan ch. apprentissage: la science}{Environ 3 paragraphes, à écrire APRÈS
le bloc précédent. C'est ici que le chapitre doit dépasser les deux articles,
qui ne se citent pas l'un l'autre et ne se savent pas parents. Trois choses à
dire: (i) la thèse du chapitre, reformulée à la lumière de ce qui précède ---
le modèle ne disparaît pas quand le critère sort, il change de rôle et devient
contrainte ou correction; (ii) le fait que ces deux rôles ne sont pas
interchangeables, et ce qui décide de l'un ou de l'autre (a-t-on une structure
à imposer, ou un biais à quantifier?); (iii) la transition vers le
\cref{ch:spdnet}, à écrire au futur: il n'y restera plus ni contrainte ni
correction, mais seulement la géométrie --- et le choix d'une moyenne y
redeviendra, par un détour, un choix de modèle statistique.}

\joinedpapers{%
  Les deux travaux dont rend compte ce chapitre sont reproduits ci-après dans
  leur version publiée.
  \joinedpaper[ECML23]{\citeauthor{ECML23}, \citetitle{ECML23}}{papers/ECML23_graphical_factor_models.pdf}
  \joinedpaper[ICML24]{\citeauthor{ICML24}, \citetitle{ICML24}}{papers/ICML24_rmt_frechet_mean.pdf}
}
